
\section{Equations and Theoretical Analysis}
\subsection{Double null foliation and equations}
In this section, we introduce the equations for Lorentzian metric under double null foliation. Let $(\mathcal{M},g)$ be a $4$-dimensional Lorentzian manifold. With double null coordinates $(u,v,\theta^A)$, the metric can be written as 
\begin{equation}
    g=-2\Omega^2(du\otimes dv+dv\otimes du)+g_{AB}(d\theta^A-b^Adu)\otimes (d\theta^B-b^Bdu),
\end{equation}
and we use $(g_{AB},b^A,\Omega)$ to represent the metric. The level sets of $u$ and $v$ are denoted by $H_u$ and $\Hb_{v}$ respectively, and their spherical intersection $H_u\cap\Hb_v$ is denoted by $S_{u,v}$.
With double null frame 
\[e_3=\Omega^{-1}(\partial_u+b),\ e_4=\Omega^{-1}\partial_v,\ e_A=\partial_{\theta^A},\]
we can define Ricci coefficients, or connection components
\begin{equation}\label{Ricci_coefficients_definition}
    \begin{aligned}
    \chi_{A B}=g\left(D_A e_4, e_B\right), &\qquad \underline{\chi}_{A B}=g\left(D_A e_3, e_B\right), \\
    \eta_A=-\frac{1}{2} g\left(D_3 e_A, e_4\right), &\qquad \underline{\eta}_A=-\frac{1}{2} g\left(D_4 e_A, e_3\right), \\
    {\omega}=-\frac{1}{4} g\left(D_4 e_3, e_4\right), &\qquad \underline{\omega}=-\frac{1}{4} g\left(D_3 e_4, e_3\right),\\
    \zeta_A=\frac{1}{2}g\left(D_A e_4,e_3\right)&=-\frac{1}{4}\Omega^{-2}(\Lie_vb^B)g_{AB}.
    \end{aligned}
\end{equation}
Their transport equations are listed below:
\begin{equation}\label{eq:Ricci_coefficients_propagation_1}
    \begin{aligned}
        \nabla_4 \operatorname{tr} \chi+\frac{1}{2}(\operatorname{tr} \chi)^2=&-\operatorname{Ric}_{44}-|\hat{\chi}|^2-2 \omega \operatorname{tr} \chi, \\
     \nabla_3 \operatorname{tr} \underline{\chi}+\frac{1}{2}(\operatorname{tr} \underline{\chi})^2=&-\operatorname{Ric}_{33}-|\underline{\hat{\chi}}|^2-2 \underline{\omega} \operatorname{tr} \underline{\chi},
    \end{aligned}
\end{equation}
\begin{equation}\label{eq:Ricci_coefficients_propagation_2}
    \begin{aligned}
        \nabla_3 \hat{\chi}_{A B}+\frac{1}{2} \operatorname{tr} \underline{\chi} \hat{\chi}_{A B}=&\widehat{\operatorname{Ric}}_{A B}+2 \underline{\omega} \hat{\chi}_{A B}+(\snab  \hat{\otimes} \eta)_{A B}+(\eta \hat{\otimes} \eta)_{A B}-\frac{1}{2} \operatorname{tr} \chi \underline{\hat{\chi}}_{A B}, \\
     \nabla_4 \underline{\hat{\chi}}_{A B}+\frac{1}{2} \operatorname{tr} \chi \underline{\hat{\chi}}_{A B}=&\widehat{\operatorname{Ric}}_{A B}+2 \omega \underline{\hat{\chi}}_{A B}+(\snab\hat{\otimes} \underline{\eta})_{A B}+(\underline{\eta} {\hat{\otimes}} \underline{\eta})_{A B}-\frac{1}{2} \operatorname{tr} \underline{\chi} \hat{\chi}_{A B},
    \end{aligned}
\end{equation}
\begin{equation}
    \begin{aligned}
        \nabla_3 \operatorname{tr} \chi=&-\frac{1}{2} \operatorname{tr} \chi \operatorname{tr} \underline{\chi}+\sg^{A B} R_{3 A B 4}+2 \underline{\omega} \operatorname{tr} \chi+2 \sdiv \eta+2|\eta|^2-\hat{\chi} \cdot \underline{\hat{\chi}}\\
        =&-\tr\chi\tr\chib+2\omegab\tr\chi+2\sdiv\eta+2\left|\eta\right|^2-2K+R+\Rc_{34}, \\
        \nabla_4 \operatorname{tr} \underline{\chi}=&-\frac{1}{2} \operatorname{tr} \chi \operatorname{tr} \underline{\chi}+\sg^{A B} R_{3 A B 4}+2 \omega \operatorname{tr} \underline{\chi}+2 \sdiv \underline{\eta}+2|\underline{\eta}|^2-\hat{\chi} \cdot \underline{\hat{\chi}}\\
        =&-\tr\chi\tr\chib+2\omega\tr\chib+2\sdiv\etab+2\left|\etab\right|^2-2K+R+\Rc_{34}.
    \end{aligned}
\end{equation} 
\begin{equation}\label{eq:Ricci_coefficients_propagation_3}
    \begin{aligned}
        \nabla_4 \eta=&-\chi \cdot(\eta-\underline{\eta})-\frac{1}{2} R_{A 434}, \\
        \nabla_3 \underline{\eta}=&-\underline{\chi} \cdot(\underline{\eta}-\eta)-\frac{1}{2} R_{A 343}, \\
        \nabla_4 \underline{\omega}=&\frac{1}{4}\left(\operatorname{Ric}_{34}+\sg^{A B} R_{3 A B 4}\right)+2 \underline{\omega} \omega+\frac{1}{2}|\eta|^2-\eta \cdot \underline{\eta}\\
        =&2\omega\omegab-\frac{1}{2}  K+\frac{1}{2} \Rc_{34}+\frac{1}{4} R +\frac{1}{4}\chih\cdot\chibh-\frac{1}{8}\tr\chi\,\tr\chib+\frac{1}{2}\left|\eta\right|^2-\eta\cdot\etab, \\
        \nabla_3 \omega=&\frac{1}{4}\left(\operatorname{Ric}_{34}+\sg^{A B} R_{3 A B 4}\right)+2 \underline{\omega} \omega+\frac{1}{2}|\underline{\eta}|^2-\eta \cdot \underline{\eta}\\
        =&2\omega\omegab-\frac{1}{2}  K+\frac{1}{2} \Rc_{34}+\frac{1}{4} R +\frac{1}{4}\chih\cdot\chibh-\frac{1}{8}\tr\chi\,\tr\chib+\frac{1}{2}\left|\underline\eta\right|^2-\eta\cdot\etab.
    \end{aligned}
\end{equation}
Moreover, we have equations for $\zeta$:
\begin{equation}\label{eq:Ricci_coefficients_propagation_4}
    \begin{aligned}
        &\Omega\Rc_{3A}=\Omega\nabla_3\zeta+\frac{3}{2}\Omega\tr\chib\zeta+\Omega\chibh\cdot\zeta+\Omega\tr\chib\nabla\log\Omega\\
        &\qquad\qquad\qquad+2\nabla(\Omega\omegab)+\dv(\Omega\chibh)-\frac{1}{2}\nabla(\Omega\tr\chib),\\
        &\Omega\Rc_{4A}=-\Omega\nabla_4\zeta-\frac{3}{2}\Omega\tr\chi\zeta-\Omega\chih\cdot\zeta+\Omega\tr\chi\nabla\log\Omega\\
        &\qquad\qquad\qquad+2\nabla(\Omega\omega)+\dv(\Omega\chih)-\frac{1}{2}\nabla(\Omega\tr\chi).
    \end{aligned}
\end{equation}
Here $K$ stands for Gaussian curvature of topological sphere $S_{u,v}$. Gauss-Codazzi equations imply 
\begin{equation}\label{eq:Gauss-Codazzi_1}
        \begin{aligned}
         K=&\frac{1}{2} \sg^{A B} R_{3 A 4 B}+\frac{1}{2} R+\frac{1}{2} \operatorname{Ric}_{34}+\frac{1}{2} \hat{\chi} \cdot \underline{\hat{\chi}}-\frac{1}{4} \operatorname{tr} \chi \operatorname{tr} \underline{\chi}, 
        \end{aligned}
        \end{equation}
\begin{equation}\label{eq:Gauss-Codazzi_2}
            \begin{aligned} 
            \snab^B \hat{\chi}_{A B}-\frac{1}{2} \snab_A \operatorname{tr} \chi=&-\frac{1}{2} R_{A 434}+\operatorname{Ric}_{4 A}+\frac{1}{2} \operatorname{tr} \chi \zeta_A-\zeta^B \hat{\chi}_{A B}, \\
            \snab^B \underline{\hat{\chi}}_{A B}-\frac{1}{2} \snab_A \operatorname{tr} \underline{\chi}=&-\frac{1}{2} R_{A 343}+\operatorname{Ric}_{3 A}-\frac{1}{2} \operatorname{tr} \underline{\chi} \zeta_A+\zeta^B \underline{\hat{\chi}}_{A B}.
            \end{aligned}
        \end{equation}
\begin{equation}\label{eq:curl-eta-etab}
            \begin{aligned}
                \snab_A \eta_B-\snab_B \eta_A=&R_{4[A B] 3}+\underline{\hat{\chi}}^C{ }_{[A} \hat{\chi}_{B] C}, \\
         \snab_A \underline{\eta}_B-\snab_B \underline{\eta}_A=&-R_{4[A B] 3}-\underline{\hat{\chi}}^C{ }_{[A} \hat{\chi}_{B] C}, \\
            \end{aligned}
        \end{equation}
For Weyl curvature,
\begin{equation*}
    W_{\mu\nu\theta\lambda}=R_{\mu\nu\theta\lambda}-\frac{1}{2}\left(\Rc_{\mu\theta}g_{\nu\lambda}+g_{\mu\theta}\Rc_{\nu\lambda}-\Rc_{\nu\theta}g_{\mu\lambda}-g_{\nu\theta}\Rc_{\mu\lambda}\right)+\frac{R}{6}\left(g_{\mu\theta}g_{\nu\lambda}-g_{\nu\theta}g_{\mu\lambda}\right),
\end{equation*}
we define the curvature components
\begin{equation}
\begin{aligned}
\beta_A & =\frac{1}{2} W\left(e_A, e_4, e_3, e_4\right), & & \underline{\beta}_A=\frac{1}{2} W\left(e_A, e_3, e_3, e_4\right), \\
\rho & =\frac{1}{4} W\left(e_4, e_3, e_4, e_3\right), & & \sigma=\frac{1}{4}{ }^* W\left(e_4, e_3, e_4, e_3\right).
\end{aligned}
\end{equation}
Then previous equations imply the following algebraic relations,
\begin{equation}
    K=-\rho+\frac{1}{2}R+\Rc_{34}+\frac{1}{2}\chibh\cdot\chih-\frac{1}{4}\tr\chib\tr\chi,
\end{equation}
\begin{equation}
    \cl\eta=-\cl\etab=\sigma+\frac{1}{2}\chibh\wedge\chih,
\end{equation}
\begin{equation}\label{Codazzi_eq_chi}
    \dv\chih-\frac{1}{2}\nabla\tr\chi=-\beta+\frac{1}{2}\Rc_{4\cdot}+\frac{1}{2}\tr\chi\zeta-\zeta\cdot\chih,
\end{equation}
\begin{equation}\label{Codazzi_eq_chib}
    \dv\chibh-\frac{1}{2}\nabla\tr\chib=\betab+\frac{1}{2}\Rc_{3\cdot}-\frac{1}{2}\tr\chib\zeta+\zeta\cdot\chibh.
\end{equation}
After introducing the reduced curvature
\begin{equation}
    \beta^r_A=\beta_A-\frac{1}{2}\Rc_{4A},\, \betab^r_A=\betab_A+\frac{1}{2}\Rc_{3A}, \,\sigma^r=\sigma+\frac{1}{2}\chibh\wedge\chih,
\end{equation}
we can write the equations for curvature components
\begin{equation}
    \begin{aligned}
        &\Omega\nabla_3 K+\Omega\tr\chib K=\dv \left(\Omega\betab^r\right)+\frac{1}{2}\dv\left(\eta\Omega\tr\chib\right)-\dv\left(\eta\cdot\Omega\chibh\right),\\
         &\Omega\nabla_4 K+\Omega\tr\chi K=-\dv \left(\Omega\beta^r\right)+\frac{1}{2}\dv\left(\etab\Omega\tr\chi\right)-\dv\left(\etab\cdot\Omega\chih\right),
    \end{aligned}
\end{equation}
\begin{equation}
    \begin{aligned}
        \nabla_3\beta^r&+\left(\tr\chib-2\omegab\right) \beta^r=-\nabla K+{}^*\nabla\sigma+2\chih\cdot\betab\\
        &\qquad+\frac{1}{2}\left(\nabla(\chih\cdot\chibh)-{}^*\nabla(\chih\wedge\chibh)\right)-\frac{1}{4}\nabla\left(\tr\chi\tr\chib\right)\\
        &\qquad+3(\eta\rho+{}^*\eta\sigma)+\nabla_A\left(g^{CD}\Rc_{CD}\right)-\nabla^B \Rc_{BA}-\frac{1}{2}\tr\chib \Rc_{4A},\\
        \nabla_4\betab^r&+\left(\tr\chi -2\omega\right)\betab^r=\nabla K+{}^*\nabla\sigma+2\chibh\cdot\beta\\
        &\qquad-\frac{1}{2}\left(\nabla(\chih\cdot\chibh)+{}^*\nabla(\chih\wedge\chibh)\right)+\frac{1}{4}\nabla\left(\tr\chi\tr\chib\right)\\
        &\qquad-3\left(\etab\rho-{}^*\etab\sigma\right) -\nabla_A(g^{CD}\Rc_{CD})+\nabla^B \Rc_{BA}+\frac{1}{2}\tr\chi \Rc_{3A}.
    \end{aligned}
\end{equation}
Using \eqref{eq:Ricci_coefficients_propagation_3}, we derive 
\begin{equation}
  \begin{aligned}
      \Omega\nabla_3\cl\etab+\Omega\tr\chib\cl\etab=&\cl(\Omega\betab^r)+\epsilon^{AC}\nabla_A\left(\Omega\chib\cdot\eta\right)_C-\epsilon^{AC}\nabla_A(\Omega\Rc)_{C3},\\
      \Omega\nabla_4\cl\eta+\Omega\tr\chi\cl\eta=&-\cl(\Omega\beta^r)+\epsilon^{AC}\nabla_A\left(\Omega\chi\cdot\etab\right)_C-\epsilon^{AC}\nabla_A(\Omega\Rc)_{C4},
  \end{aligned}
\end{equation}
and
\begin{equation}
    \begin{aligned}
        &\Omega\nabla_3\left(\dv \etab-K\right)+\Omega\tr\chib\left(\dv \etab-K\right)=4\dv(\Omega\chibh\cdot\zeta)-\dv(\Omega\Rc_3),\\
        &\Omega\nabla_4\left(\dv \eta-K\right)+\Omega\tr\chi\left(\dv \eta-K\right)=-4\dv(\Omega\chih\cdot\zeta)-\dv(\Omega\Rc_4).
    \end{aligned}
\end{equation}
For a scalar function $\phi$, there is equation
\begin{equation}\label{eq:ese-wave-eq}
    -\Omega^2\square\phi=\Omega e_3(\Omega e_4\phi)+\frac{1}{2}\Omega\tr\chi\Omega e_3\phi+\frac{1}{2}\Omega\tr\chib\Omega e_4\phi-\Omega^2\Delta\phi-2\Omega^2\eta\cdot\nabla\phi.
\end{equation}

\subsection{Local existence}
We now introduce some results on Einstein vacuum equations. Similar results can be derived for other Einstein field equations and we omit them due to length of the article.

We first state the existence theorem for characteristic initial value problem of Einstein vacuum equations from \cite{luk-12}.
\begin{theorem}
    Given initial null hypersurface $H_0,\Hb_0$ and regular initial metric $(g,\Omega,b)$ on it, we define $\chi_{AB}=\frac{1}{2}\Lie_{e_4}g_{AB}$ on $H_0$ and $\chib_{AB}=\frac{1}{2}\Lie_{e_3}g_{AB}$ on $\Hb_0$. Suppose 
    \[e_4(\tr\chi)+\frac{1}{2}(\tr\chi)^2+|\chih|^2+2\omega\tr\chi=0,\ e_3(\tr\chib)+\frac{1}{2}(\tr\chib)^2+|\chibh|^2+2\omegab\tr\chib,\]
    hold on $H_0,\Hb_0$ respectively, then there is a future open neighborhood of $H_0\cup \Hb_0$ and metric $(g,b,\Omega)$ on it solving the Einstein vacuum equations and satisfying the initial data.
\end{theorem}
This theorem reveals the local existence, and we further remark that for larger region, such as $(0,1)^2\times\mathbb{S}^2$, as long as the a priori estimates hold in it, the solution can be extended to the region.

\subsection{Residual control}
In this paper, we target to deal with asymmetric Einstein equations, which in general do not have explicit solutions. To measure the computational error, we use Ricci residues. For Einstein equation 
\[G_{\mu\nu}:=\Rc_{\mu\nu}-\frac{1}{2}Rg_{\mu\nu}=T_{\mu\nu},\]
we measure $\|G-T\|_{L^2(S_{u,v})}$ to evaluate the approximation. The theorems below reveals that in our case the difference of metrics is controlled by the difference of curvature residues.


\begin{theorem}
    We consider metric $\oc{g}_{AB},\oc{b},\oc{\Omega}$ on $(0, 1)\times (0,1)\times \mathbb{S}^2$ such that
    \[\oc{\Rc}_{44}(0,v)=0,\ \oc{\Rc}_{33}(u,0)=0,\ \oc{\Rc}_{4A}(0,v)=0,\ \oc{\Rc}_{3A}(u,0)=0,\]
    % \[\oc{\Rc}_{AB}(u,0)=0,\ \oc{\Rc_{34}}(u,0)=0,\ \oc{R}(u,0)=0,\]
    and estimates 
    \begin{equation}
        \|\oc{\Rc}_{AB},\oc{\Rc}_3,\oc{\Rc}_4,\oc{\Rc}_{34},\oc{\Rc}_{33},\oc{\Rc}_{44}\|_{L^\infty_uL^\infty_vH^3(S)}\leq \epsilon.
    \end{equation}
    We assume moreover 
    \begin{equation}
        \| \partial_{\mathbb{S}}\oc{g}, \partial_{\mathbb{S}}\oc{b}, \partial_{\mathbb{S}}\oc\Omega, \oc{\psi}\|_{L^\infty_u L^\infty_v H^4(S)} \leq C_0,
    \end{equation}
    where $\psi$ stands for Ricci coefficients in interest.
    Then for $\epsilon=\epsilon(C_0)$ sufficiently small, there exists $(g,b,\Omega)$ solving EVE on $( 0,1)\times (0,1)\times \mathbb{S}^2$ with initial data 
    \[(g,b,\Omega)=(\oc{g},\oc{b},\oc\Omega),\ {\rm on}\ \Hb_0\cup H_{0},\]
    and estimates
    \begin{equation}
        \|g-\oc{g},b-\oc{b},\Omega-\oc{\Omega}\|_{L^\infty_u L^\infty_v H^4(S)}\leq \epsilon^{1/2}.
    \end{equation}
\end{theorem}

The methods of deriving the equations of difference and doing estimates are similar to section 7 of \cite{an-26},
so we only write the proof in sketch. 

Define auxiliary functions $\omega^\dagger,\omegab^\dagger$ by
\[\nabla_3\omega^\dagger=\frac{1}{2}\cl\zeta,\ \omega^\dagger|_{H_{-1}}=0,\ \nabla_4\omegab^\dagger=\frac{1}{2}\cl\zeta,\ \omegab^\dagger|_{\Hb_0}=0.\]
We make bootstrap assumption
\begin{equation*}
    \begin{aligned}
        &\|\wt{\beta},\wt{K},\nabla\wt\eta,\nabla\wt{\etab},\nabla\wt{\chi},\nabla\wt{\omega},\nabla\wt{\omega^\dagger}\|_{L^2_vH^3(S)}\\
        &\qquad+\|\wt{K},\wt{\betab},\nabla\wt{\eta},\nabla\wt\etab,\nabla\wt{\chib},\nabla\wt{\omegab},\nabla\wt{\omegab^\dagger}\|_{L^2_uH^3(S)}\leq B\exp(D(v+u))\epsilon.
    \end{aligned}
\end{equation*}
We will prove that we can choose $D$ sufficiently large so that the bootstrap constant $B$ can be improved to $B\leq 1$.

The motivation is direct. Let $\partial_uf\leq C_0f,\partial_vf\leq C_0f$, and suppose
$$f(u,v)\leq B\exp(D(v+u))\epsilon,$$ then we integrate it 
\begin{equation*}
    \begin{aligned}
        f(u,v)-f(0,v)\leq&C_0\int_{0}^u f(u^\prime,v) du^\prime\\
        \leq  & \frac{C_0B}{D}\left(\exp(D(u+v))-\exp(Dv)\right)\epsilon \leq \frac{C_0B}{D} \exp(D(u+v))\epsilon,\\
        f(u,v)-f(u,0)\leq & C_0\int_{0}^v f(u,v^\prime) dv^\prime\\
        \leq & \frac{C_0B}{D}\left(\exp(D(u+v))-\exp(Du)\right)\epsilon \leq \frac{C_0B}{D} \exp(D(u+v))\epsilon.
    \end{aligned}
\end{equation*}
If we choose $D^{1/2}\geq C_0B$, the bootstrap assumption will be closed. Absorbing the bootstrap constant by the parameter $D$ actually shows the increasing rate is controlled exponentially. Even though the norms increase quickly, they will not exceed $O(\epsilon)\leq o(\epsilon^{1/2})$ in finite region.

In summary, we conclude 
\begin{lemma}
    Suppose $\|f\|_{L^2(S_{u,v})}\leq B\exp(D(v+u))\epsilon$, then we have
    \begin{equation}
        \|f\|_{L^2_uL^2(S)}+\|f\|_{L^2_vL^2(S)} \leq D^{-1/2}\exp(D(v+u))\epsilon,
    \end{equation}
    when $D$ is sufficiently large.

    Suppose either $\|h\|_{L^2_uL^2(S)}\leq B\exp(D(u+v))\epsilon$ or $\|h\|_{L^2_vL^2(S)}\leq B\exp(D(u+v))\epsilon$, then we have    
    \begin{equation}
        \|h\|_{L^2_uL^2_vL^2(S)} \leq D^{-1/2}\exp(D(v+u))\epsilon,
    \end{equation}
    when $D$ is sufficiently large.
\end{lemma}

We use notation
\[\psi_1\in\{\eta,\etab,\chi,\omega\},\ \psi_2\in\{\eta,\etab,\chib,\omegab\}\]
Because 
\[\Lie_v b=-2\Omega^2(\eta-\etab),\ \Lie_vg=2\Omega\chi,\ \Lie_v\log\Omega=-2\Omega\omega,\]
we obtain that
\begin{equation}
    \|\wt{b},\wt{g},\wt{\log\Omega}\|_{H^4(S_{u,v})}\lesssim \|\wt{\eta},\wt{\etab},\wt{\chi},\wt{\omega},\wt{\Omega},\wt{g}\|_{L^2_vH^4(S)}\lesssim B\exp(D(u+v))\epsilon.
\end{equation}
For Bianchi pair $(\Psi_1,\Psi_2)$ in interest, we consider equations 
\begin{equation*}
    \begin{aligned}
        \Omega\nabla_3\wt{\Psi_1} &+\mathcal{D}\wt{\Psi_2}=\wt{\psi\Psi}+\wt{\psi\nabla\psi}+\wt{\psi^3}+\wt{b}\cdot\nabla\oc{\Psi_1}+\wt{\partial g}\oc{\Psi_2}+\wt{g}\oc{\nabla\Psi_2}+\oc\psi\oc\Rc+\oc{\nabla\Rc},\\
        \Omega\nabla_4\wt{\Psi_2} &-{}^*\mathcal{D}\wt{\Psi_1}=\wt{\psi\Psi}+\wt{\psi\nabla\psi}+\wt{\psi^3}+\wt{\partial g}\oc{\Psi_1}+\wt{g}\oc{\nabla\Psi_1}+\oc\psi\oc\Rc+\oc{\nabla\Rc}.
    \end{aligned}
\end{equation*}
Energy estimates give that 
\begin{equation*}
    \begin{aligned}
        &\|\wt\Psi_1\|^2_{L^2_vH^3(S)}+\|\wt\Psi_2\|^2_{L^2_uH^3(S)}\\
        &\qquad\qquad\lesssim \epsilon^2+\|\wt{\partial g},\wt{g},\wt{\psi},\wt{\nabla\psi},\wt\Psi,\wt{b}\|^2_{L^2_uL^2_vH^3(S)}\lesssim (1+D^{-1/2})\exp(2D(u+v))\epsilon^2.
    \end{aligned}
\end{equation*}
Applying the estimate to pairs
\[(\beta,(K,\cl\eta)),\ ((K,\cl\etab),\betab),\ (\nabla\tr\chib,0),\ (0,\nabla\tr\chi),\]
\[ (\dv\etab-K,0),\ (0,\dv\eta-K),\ (\nabla\omega+{}^*\nabla\omega^\dagger-\frac{1}{2}\beta,0),\ (0,-\nabla\omegab+{}^*\nabla\omegab^\dagger-\frac{1}{2}\betab),\]
will finish the proof. 
\begin{remark}
    The conditions 
    \[\oc{\Rc}_{44}(-1,v)=0,\ \oc{\Rc}_{33}(u,0)=0,\]
    guarantee that $(\oc{g},\oc{b},\oc{\Omega})$ is a valid initial data set for characteristic initial value problem of Einstein vacuum equation, and 
    \[\oc{\Rc}_{4A}(-1,v)=0,\ \oc{\Rc}_{3A}(u,0)=0,\]
    ensure that $\oc\zeta$ along $H_{-1}$ and $\Hb_0$ is equal to the $\zeta$ of corresponding EVE solution. 
\end{remark}

